\documentclass[border=5pt]{standalone}
\usepackage{amsmath}
\usepackage{tikz}
\usepackage{tikz-3dplot}
\usetikzlibrary{arrows.meta}
\usetikzlibrary{fadings}

% Radial fading used for the blade "pearl" specular highlight: opaque at
% the circle center, fully transparent at the circle's edge. Renders as a
% soft round shine spot when applied as `path fading=glossspot` to a
% filled circle.
\tikzfading[name=glossspot,
  inner color=transparent!0,
  outer color=transparent!100]

% --- Local macro shims (the standalone document does not load the thesis preamble)
\newcommand{\Tgroup}{\mathcal{T}}
\newcommand{\Ogroup}{\mathcal{O}}
\newcommand{\Igroup}{\mathcal{I}}
\newcommand{\TwoT}{2\Tgroup}
\newcommand{\TwoO}{2\Ogroup}
\newcommand{\TwoI}{2\Igroup}
\newcommand{\phig}{\tau}
\newcommand{\invphig}{\sigma}
\newcommand{\ket}[1]{|#1\rangle}

% --- View setup (matches platonic_solids.tex) -----------------------------
\tdplotsetmaincoords{60}{120}
\pgfmathsetmacro{\viewX}{sin(60)*sin(120)}
\pgfmathsetmacro{\viewY}{-sin(60)*cos(120)}
\pgfmathsetmacro{\viewZ}{cos(60)}

% --- Geometry constants ---------------------------------------------------
\pgfmathsetmacro{\PHI}{(1+sqrt(5))/2}
\pgfmathsetmacro{\IPHI}{(sqrt(5)-1)/2}
\pgfmathsetmacro{\TET}{1/sqrt(3)}
\pgfmathsetmacro{\ICOS}{1/sqrt(2+\PHI)}
\pgfmathsetmacro{\PHIICOS}{\PHI*\ICOS}
\pgfmathsetmacro{\PHITET}{\PHI*\TET}
\pgfmathsetmacro{\IPHITET}{\IPHI*\TET}

% --- Vertex storage (matches platonic_solids.tex) -------------------------
\newcommand{\PVert}[5]{%
  \expandafter\edef\csname Vx#1at#2\endcsname{#3}%
  \expandafter\edef\csname Vy#1at#2\endcsname{#4}%
  \expandafter\edef\csname Vz#1at#2\endcsname{#5}%
}
\newcommand{\Vx}[2]{\csname Vx#1at#2\endcsname}
\newcommand{\Vy}[2]{\csname Vy#1at#2\endcsname}
\newcommand{\Vz}[2]{\csname Vz#1at#2\endcsname}

% --- Vertex tables (copied verbatim from platonic_solids.tex) -------------

% Tetrahedron: cube-inscribed, even sign parity, scaled 1/sqrt(3)
\PVert{tet}{1}{\TET}{\TET}{\TET}
\PVert{tet}{2}{-\TET}{-\TET}{\TET}
\PVert{tet}{3}{-\TET}{\TET}{-\TET}
\PVert{tet}{4}{\TET}{-\TET}{-\TET}

% Cube: (+/-, +/-, +/-)/sqrt(3)
\PVert{cube}{1}{-\TET}{-\TET}{-\TET}
\PVert{cube}{2}{\TET}{-\TET}{-\TET}
\PVert{cube}{3}{\TET}{\TET}{-\TET}
\PVert{cube}{4}{-\TET}{\TET}{-\TET}
\PVert{cube}{5}{-\TET}{-\TET}{\TET}
\PVert{cube}{6}{\TET}{-\TET}{\TET}
\PVert{cube}{7}{\TET}{\TET}{\TET}
\PVert{cube}{8}{-\TET}{\TET}{\TET}

% Octahedron: +-e_i
\PVert{oct}{1}{1}{0}{0}
\PVert{oct}{2}{-1}{0}{0}
\PVert{oct}{3}{0}{1}{0}
\PVert{oct}{4}{0}{-1}{0}
\PVert{oct}{5}{0}{0}{1}
\PVert{oct}{6}{0}{0}{-1}

% Icosahedron: cyclic perms of (0, +-phi, +-1), scaled 1/sqrt(1+phi^2)
\PVert{ico}{1}{\PHIICOS}{\ICOS}{0}
\PVert{ico}{2}{-\PHIICOS}{\ICOS}{0}
\PVert{ico}{3}{\PHIICOS}{-\ICOS}{0}
\PVert{ico}{4}{-\PHIICOS}{-\ICOS}{0}
\PVert{ico}{5}{0}{\PHIICOS}{\ICOS}
\PVert{ico}{6}{0}{-\PHIICOS}{\ICOS}
\PVert{ico}{7}{0}{\PHIICOS}{-\ICOS}
\PVert{ico}{8}{0}{-\PHIICOS}{-\ICOS}
\PVert{ico}{9}{\ICOS}{0}{\PHIICOS}
\PVert{ico}{10}{-\ICOS}{0}{\PHIICOS}
\PVert{ico}{11}{\ICOS}{0}{-\PHIICOS}
\PVert{ico}{12}{-\ICOS}{0}{-\PHIICOS}

% Dodecahedron: 8 cube vertices + cyclic perms of (0, +-1/phi, +-phi), all /sqrt(3)
\PVert{dod}{1}{\TET}{\TET}{\TET}
\PVert{dod}{2}{\TET}{\TET}{-\TET}
\PVert{dod}{3}{\TET}{-\TET}{\TET}
\PVert{dod}{4}{\TET}{-\TET}{-\TET}
\PVert{dod}{5}{-\TET}{\TET}{\TET}
\PVert{dod}{6}{-\TET}{\TET}{-\TET}
\PVert{dod}{7}{-\TET}{-\TET}{\TET}
\PVert{dod}{8}{-\TET}{-\TET}{-\TET}
\PVert{dod}{9}{0}{\IPHITET}{\PHITET}
\PVert{dod}{10}{0}{\IPHITET}{-\PHITET}
\PVert{dod}{11}{0}{-\IPHITET}{\PHITET}
\PVert{dod}{12}{0}{-\IPHITET}{-\PHITET}
\PVert{dod}{13}{\IPHITET}{\PHITET}{0}
\PVert{dod}{14}{\IPHITET}{-\PHITET}{0}
\PVert{dod}{15}{-\IPHITET}{\PHITET}{0}
\PVert{dod}{16}{-\IPHITET}{-\PHITET}{0}
\PVert{dod}{17}{\PHITET}{0}{\IPHITET}
\PVert{dod}{18}{\PHITET}{0}{-\IPHITET}
\PVert{dod}{19}{-\PHITET}{0}{\IPHITET}
\PVert{dod}{20}{-\PHITET}{0}{-\IPHITET}

% --- POVM-specific helpers ------------------------------------------------

% Project a stored vertex onto the camera direction; result in \NdotV.
\newcommand{\vertexNdotV}[2]{%
  \pgfmathsetmacro{\NdotV}{%
    \Vx{#1}{#2}*\viewX + \Vy{#1}{#2}*\viewY + \Vz{#1}{#2}*\viewZ}%
}

% Edge: continuous depth gradient. Shade is computed from the average of
% the two endpoints' n.v (proxy for "how front-facing is this edge on
% average"). Mapped linearly into black!24..75, so a fully back edge reads
% as a faint receding scaffold and a fully front edge reads as a
% dark structural line, with everything in between sweeping smoothly.
% Each edge is routed to backedges or frontedges layer by sign of \NdotVmid;
% the front layer additionally takes a white preaction halo (width \haloW,
% set per panel) so that foreground edges occlude crossed background edges.
\newcommand{\drawEdge}[3]{%
  \vertexNdotV{#1}{#2}\let\NdotVa\NdotV
  \vertexNdotV{#1}{#3}\let\NdotVb\NdotV
  \pgfmathsetmacro{\NdotVmid}{(\NdotVa + \NdotVb) / 2}%
  \pgfmathsetmacro{\edgeShade}{int(round(max(24, min(75, 48 + 50*\NdotVmid))))}%
  \pgfmathparse{\NdotVmid > 0 ? 1 : 0}%
  \ifnum\pgfmathresult=1
    \begin{pgfonlayer}{frontedges}
      \draw[edge, draw=black!\edgeShade,
            preaction={draw=white, line width=\haloW, line cap=round}]
        (\Vx{#1}{#2},\Vy{#1}{#2},\Vz{#1}{#2}) --
        (\Vx{#1}{#3},\Vy{#1}{#3},\Vz{#1}{#3});
    \end{pgfonlayer}
  \else
    \begin{pgfonlayer}{backedges}
      \draw[edge, draw=black!\edgeShade,
            preaction={draw=white, line width=\haloBackW, line cap=round}]
        (\Vx{#1}{#2},\Vy{#1}{#2},\Vz{#1}{#2}) --
        (\Vx{#1}{#3},\Vy{#1}{#3},\Vz{#1}{#3});
    \end{pgfonlayer}
  \fi
}

% 3D blade arrow geometry. Each POVM ray is drawn as a 7-vertex flat
% arrow polygon living in a specific plane in 3D space, not as a 2D stroke
% that always faces the camera. Width axis w = normalize(v x U), where U
% is a per-panel reference direction (\bladeUx, \bladeUy, \bladeUz). The
% blade lies in the plane spanned by v and w. Because w is fixed in world
% space, the camera foreshortens the blade by orientation: blades whose
% width axis lies in the screen plane look wide; blades whose width axis
% points along the camera direction look edge-on. Per-panel U is tuned so
% that each panel's blades land at visually distinct widths without
% disappearing edge-on:
%   - Octahedron uses U = z_hat. This places +/-x and +/-y blades flat in
%     the xy-plane (the axes the Pauli kets live on); +/-z rays fall back
%     to w = y_hat (yz-plane).
%   - Tetrahedron / Cube use U = (0,1,1)/sqrt(2). With U = z_hat the
%     chiral (1,1,1)/sqrt(3) rays foreshorten too aggressively (one tet
%     blade drops to ~7% width). Tilting U keeps the minimum tet width
%     around 23%, so every blade reads as a blade.
%
% Shape: shaft + barbed arrowhead. Shoulder (where shaft ends and the head
% begins) is at \bladeShoulderT along v; wings (head's widest point) are
% at \bladeWingT, slightly *behind* the shoulder so the head has small
% backward-swept barbs. The head tapers from the wing inward to the tip
% at v. Shaft and arrowhead are a single closed polygon, so they merge
% seamlessly. Half-widths are in 3D units (the enclosing scope rescales
% by 2.6 for screen size).
%
% Fill: two-tone (lit + shadow) split across the centerline (the segment
% from origin to v). Right half (+w side) is lighter, left half (-w side)
% is darker, simulating a glossy blade lit from the +w direction. A crisp
% black outline binds the whole polygon. The lit half is slightly more
% translucent than the shadow half so the contrast reads as "shine" rather
% than "two flat panels".
\newcommand{\bladeShaftHalfW}{0.017}
\newcommand{\bladeHeadHalfW}{0.048}
\newcommand{\bladeShoulderT}{0.83}
\newcommand{\bladeWingT}{0.78}

% Default U = z_hat (octahedron default); panels override before drawing.
\newcommand{\bladeUx}{0}
\newcommand{\bladeUy}{0}
\newcommand{\bladeUz}{1}

\newcommand{\drawBlade}[2]{% {solidName}{vertexIdx}
  \pgfmathsetmacro{\bvx}{\Vx{#1}{#2}}%
  \pgfmathsetmacro{\bvy}{\Vy{#1}{#2}}%
  \pgfmathsetmacro{\bvz}{\Vz{#1}{#2}}%
  % w_raw = v x U. Normalize; fallback to (0,1,0) if v parallel to U.
  \pgfmathsetmacro{\bwxr}{\bvy*\bladeUz - \bvz*\bladeUy}%
  \pgfmathsetmacro{\bwyr}{\bvz*\bladeUx - \bvx*\bladeUz}%
  \pgfmathsetmacro{\bwzr}{\bvx*\bladeUy - \bvy*\bladeUx}%
  \pgfmathsetmacro{\bwn}{sqrt(\bwxr*\bwxr + \bwyr*\bwyr + \bwzr*\bwzr)}%
  \pgfmathparse{\bwn > 0.01 ? 1 : 0}%
  \ifnum\pgfmathresult=1
    \pgfmathsetmacro{\bwx}{\bwxr/\bwn}%
    \pgfmathsetmacro{\bwy}{\bwyr/\bwn}%
    \pgfmathsetmacro{\bwz}{\bwzr/\bwn}%
  \else
    \pgfmathsetmacro{\bwx}{0}%
    \pgfmathsetmacro{\bwy}{1}%
    \pgfmathsetmacro{\bwz}{0}%
  \fi
  % Shaft / head offsets along w, shoulder / wing v-positions along v.
  \pgfmathsetmacro{\bSx}{\bladeShaftHalfW*\bwx}%
  \pgfmathsetmacro{\bSy}{\bladeShaftHalfW*\bwy}%
  \pgfmathsetmacro{\bSz}{\bladeShaftHalfW*\bwz}%
  \pgfmathsetmacro{\bHx}{\bladeHeadHalfW*\bwx}%
  \pgfmathsetmacro{\bHy}{\bladeHeadHalfW*\bwy}%
  \pgfmathsetmacro{\bHz}{\bladeHeadHalfW*\bwz}%
  \pgfmathsetmacro{\bShx}{\bladeShoulderT*\bvx}%
  \pgfmathsetmacro{\bShy}{\bladeShoulderT*\bvy}%
  \pgfmathsetmacro{\bShz}{\bladeShoulderT*\bvz}%
  \pgfmathsetmacro{\bWx}{\bladeWingT*\bvx}%
  \pgfmathsetmacro{\bWy}{\bladeWingT*\bvy}%
  \pgfmathsetmacro{\bWz}{\bladeWingT*\bvz}%
  % All blades go to `bladelayer` (between backedges and frontedges) so
  % the polytope front edges with their white halos render on top.
  \begin{pgfonlayer}{bladelayer}\drawBladePoly\end{pgfonlayer}
}

% Gel-button blade: solid base + specular sheen + outline.
% Base fill covers the full polygon. The sheen is a white-to-dark gradient
% applied to the +w half polygon, with `shading angle` chosen so the bright
% end of the gradient sits at the +w outer edge in screen space. For
% tdplot's (theta=60, phi=120) view, the screen projection of a 3D
% direction (wx, wy, wz) is approximately
%   screen_x = -0.866 wx - 0.5 wy
%   screen_y = 0.25 wx - 0.433 wy + 0.866 wz,
% so atan2(screen_y, screen_x) gives the +w screen-space angle; subtracting
% 90 deg converts to a TikZ shading angle (default "top" is at +y/north).
\newcommand{\drawBladePoly}{%
  % Base fill (full polygon).
  \fill[bladebase]
    (-\bSx,-\bSy,-\bSz) --
    (\bSx,\bSy,\bSz) --
    (\bShx+\bSx,\bShy+\bSy,\bShz+\bSz) --
    (\bWx+\bHx,\bWy+\bHy,\bWz+\bHz) --
    (\bvx,\bvy,\bvz) --
    (\bWx-\bHx,\bWy-\bHy,\bWz-\bHz) --
    (\bShx-\bSx,\bShy-\bSy,\bShz-\bSz) --
    cycle;
  % Per-blade fading angle. The sheen should run perpendicular to the
  % projected ray direction in screen, with its opaque end on the side
  % where +w lands. We can't just use the projected +w direction, since
  % 3D perpendicularity (v perp w) doesn't survive the projection -- for
  % thin foreshortened blades that puts the gradient diagonally across the
  % long axis. Instead, project v, take its screen-perpendicular, and pick
  % the perpendicular's sign from the projection of w.
  \pgfmathsetmacro{\bvSx}{-0.866*\bvx - 0.5*\bvy}%
  \pgfmathsetmacro{\bvSy}{0.25*\bvx - 0.433*\bvy + 0.866*\bvz}%
  \pgfmathsetmacro{\bwSx}{-0.866*\bwx - 0.5*\bwy}%
  \pgfmathsetmacro{\bwSy}{0.25*\bwx - 0.433*\bwy + 0.866*\bwz}%
  % Does +w project to the (-vSy, vSx) side of v in screen?
  \pgfmathsetmacro{\bSheenSign}{(\bwSy*\bvSx - \bwSx*\bvSy) > 0 ? 1 : -1}%
  \pgfmathsetmacro{\bSheenAng}{atan2(\bSheenSign*\bvSx, -\bSheenSign*\bvSy) - 90}%
  % Soft sheen overlay on +w half polygon: white fill that fades from
  % opaque (at the +w outer edge) to fully transparent (at the centerline),
  % so the centerline blends seamlessly into the dark base. This is the
  % broad "ambient lit half" component.
  \fill[bladesheenSoft, fading angle=\bSheenAng]
    (\bSx,\bSy,\bSz) --
    (\bShx+\bSx,\bShy+\bSy,\bShz+\bSz) --
    (\bWx+\bHx,\bWy+\bHy,\bWz+\bHz) --
    (\bvx,\bvy,\bvz) --
    (0,0,0) --
    cycle;
  % Outline first, so the stripe (drawn after) sits on top of the +w
  % outer edge instead of being clipped by the dark stroke.
  \draw[bladeoutline]
    (-\bSx,-\bSy,-\bSz) --
    (\bSx,\bSy,\bSz) --
    (\bShx+\bSx,\bShy+\bSy,\bShz+\bSz) --
    (\bWx+\bHx,\bWy+\bHy,\bWz+\bHz) --
    (\bvx,\bvy,\bvz) --
    (\bWx-\bHx,\bWy-\bHy,\bWz-\bHz) --
    (\bShx-\bSx,\bShy-\bSy,\bShz-\bSz) --
    cycle;
  % Pearl specular highlight: a single radial shine spot anchored at the
  % shoulder-right corner of the +w outer edge (where the shaft meets the
  % head's barb -- visually the most distinctive point of the lit silhouette).
  % The circle is filled with the `glossspot` radial fading so the spot is
  % opaque at center and smoothly fades to transparent at the perimeter,
  % then clipped to the full polygon so any overflow past the blade's edges
  % is trimmed cleanly. Drawn AFTER the outline so it sits on top.
  \begin{scope}
    \clip
      (-\bSx,-\bSy,-\bSz) --
      (\bSx,\bSy,\bSz) --
      (\bShx+\bSx,\bShy+\bSy,\bShz+\bSz) --
      (\bWx+\bHx,\bWy+\bHy,\bWz+\bHz) --
      (\bvx,\bvy,\bvz) --
      (\bWx-\bHx,\bWy-\bHy,\bWz-\bHz) --
      (\bShx-\bSx,\bShy-\bSy,\bShz-\bSz) --
      cycle;
    \fill[bladesheenPearl]
      (\bShx+\bSx,\bShy+\bSy,\bShz+\bSz) circle (2.6mm);
  \end{scope}
}

% Bloch ray (3D blade) + vertex disk. Blade depth-routed via \drawBlade;
% disks keep their original front/back treatment (back disks smaller at
% 0.85x, lighter). Disk drawn after the blade so it caps the tip.
\newcommand{\drawRayDisk}[3]{% {solidName}{vertexIdx}{diskRadius_mm}
  \drawBlade{#1}{#2}%
  \vertexNdotV{#1}{#2}%
  \pgfmathparse{\NdotV>0 ? 1 : 0}%
  \ifnum\pgfmathresult=1
    \fill[diskfront]
      (\Vx{#1}{#2},\Vy{#1}{#2},\Vz{#1}{#2}) circle (#3mm);
  \else
    \pgfmathsetmacro{\backR}{#3*0.85}%
    \fill[diskback]
      (\Vx{#1}{#2},\Vy{#1}{#2},\Vz{#1}{#2}) circle (\backR mm);
  \fi
}

% Disk only (no ray). Used in the icosahedron and dodecahedron panels,
% where 12-20 rays would starburst across the polytope interior.
\newcommand{\drawDisk}[3]{% {solidName}{vertexIdx}{diskRadius_mm}
  \vertexNdotV{#1}{#2}%
  \pgfmathparse{\NdotV>0 ? 1 : 0}%
  \ifnum\pgfmathresult=1
    \fill[diskfront]
      (\Vx{#1}{#2},\Vy{#1}{#2},\Vz{#1}{#2}) circle (#3mm);
  \else
    \pgfmathsetmacro{\backR}{#3*0.85}%
    \fill[diskback]
      (\Vx{#1}{#2},\Vy{#1}{#2},\Vz{#1}{#2}) circle (\backR mm);
  \fi
}

% Six Pauli-eigenstate labels on the unit sphere. Anchors push each label outward
% along its screen-projected axis direction (computed for the {60}{120} view), so
% labels sit just outside the silhouette where possible. Drawn last so they sit on top.
\newcommand{\PauliLabels}{%
  \node[font=\scriptsize, anchor=south]      at (0,0,1)  {$\ket{0}$};
  \node[font=\scriptsize, anchor=north]      at (0,0,-1) {$\ket{1}$};
  \node[font=\scriptsize, anchor=north east] at (1,0,0)  {$\ket{+}$};
  \node[font=\scriptsize, anchor=south west] at (-1,0,0) {$\ket{-}$};
  \node[font=\scriptsize, anchor=west]       at (0,1,0)  {$\ket{{+}i}$};
  \node[font=\scriptsize, anchor=east]       at (0,-1,0) {$\ket{{-}i}$};
}

% Short radial nubs at the front piercings of +x, +y, +z. Used in panels where
% the Pauli-eigenstate kets read as vertex labels and would create ambiguity;
% nubs supply orientation without text.
\newcommand{\AxisNubs}{%
  \draw[axisnub] (1,0,0) -- (1.06,0,0);
  \draw[axisnub] (0,1,0) -- (0,1.06,0);
  \draw[axisnub] (0,0,1) -- (0,0,1.06);
}

% --- Edge layers + halo width --------------------------------------------
% Two declared layers split edges by depth so foreground edges always render
% after background edges, regardless of source order. The white preaction on
% the front layer (controlled by \haloW, set per panel) knocks out crossed
% background edges at intersections.
\pgfdeclarelayer{backedges}
\pgfdeclarelayer{bladelayer}
\pgfdeclarelayer{frontedges}
\pgfsetlayers{backedges,bladelayer,frontedges,main}
\newcommand{\haloW}{0pt}
\newcommand{\haloBackW}{0pt}
\newcommand{\haloFrontEqW}{0pt}

% --- TikZ styles ----------------------------------------------------------
% Style values are visual-iteration tunables. Edge shading is continuous:
% \drawEdge maps the midpoint n.v of each edge into black!35..75 so depth
% reads as a smooth gradient rather than discrete tiers (see the macro for
% the mapping). Disk fill, the 0.85 back-disk scale, and the blade opacity
% also control depth contrast.
%
% The blade style is the fill+stroke applied to the 3D POVM arrow polygons
% built by \drawBlade. Glossy-black: translucent black fill + thin black
% outline. Translucency (fill opacity 0.35) lets the wireframe behind show
% through; the outline keeps the blade's silhouette crisp at any
% foreshortening angle, even when the blade is nearly edge-on.
\tikzset{
  edge/.style           = {line width=0.4pt, line cap=round},
  bladebase/.style        = {fill=black, fill opacity=0.75},
  bladesheenSoft/.style   = {fill=white, fill opacity=0.2, path fading=north},
  bladesheenPearl/.style  = {fill=white, fill opacity=0.5, path fading=glossspot},
  bladeoutline/.style   = {draw=black, draw opacity=0.95,
                           line width=0.32pt, line join=round, line cap=round},
  axis/.style           = {draw=black!22, line width=0.3pt, line cap=round, dotted},
  axisnub/.style        = {draw=black!55, line width=0.5pt, line cap=round},
  diskfront/.style      = {fill=black},
  diskback/.style       = {fill=black!35},
  povmsphere/.style     = {draw=black!45, line width=0.35pt},
  povmsphereback/.style = {draw=black!22, line width=0.35pt},
  rowlabel/.style       = {anchor=east, font=\normalsize},
  panelname/.style      = {anchor=north, font=\footnotesize},
}

% --- Per-panel macros -----------------------------------------------------
% \PanelHeader opens an outer xshift/yshift scope, draws the screen-plane
% silhouette, and opens an inner tdplot scope. \PanelFooter closes the
% inner scope, places the name + info block, and closes the outer scope.
% Three coordinate axes through origin, drawn dotted-faint as orientation guides.
% Treated uniformly as scaffolding "behind" the polytope: negative halves go
% to backedges, positive halves to frontedges -- so polytope edges in the
% same layer overdraw the axes (drawn first within the layer via PanelHeader),
% and front-edge halos chop them at crossings. Geometrically not exact at
% every point (axes are diameters through origin, so depth along each half
% varies from 0 at origin to viewN at the unit sphere), but it's the correct
% direction of approximation: an axis interior is inside the unit sphere and
% is more often behind front polytope structure than in front of it.
\newcommand{\Axes}{%
  \begin{pgfonlayer}{backedges}
    \draw[axis] (-1,0,0) -- (0,0,0);
    \draw[axis] (0,-1,0) -- (0,0,0);
    \draw[axis] (0,0,-1) -- (0,0,0);
  \end{pgfonlayer}
  \begin{pgfonlayer}{frontedges}
    \draw[axis] (0,0,0) -- (1,0,0);
    \draw[axis] (0,0,0) -- (0,1,0);
    \draw[axis] (0,0,0) -- (0,0,1);
  \end{pgfonlayer}
}

% Equator (xy-plane great circle): solid arc for the camera-facing half,
% faded solid for the far side (no dashes -- matches the edge treatment).
% Front/back boundary at theta = -60, +120 deg for the {60}{120} view
% (camera direction (0.75, 0.433, 0.5)).
% Asymmetric depth routing: the back arc goes to backedges so polytope edges
% render over it, the front arc stays in main so it draws over polytope
% edges and is not knocked out by their white halos at crossings. The front
% arc additionally takes a white halo (\haloFrontEqW) so that polytope edges
% are cleanly broken where they cross the equator -- geometrically correct
% since the equator lies on the unit sphere and polytope chords are inside.
\newcommand{\Equator}{%
  \draw[povmsphere,
        preaction={draw=white, line width=\haloFrontEqW, line cap=round}]
    plot[domain=-60:120, samples=37, variable=\t, smooth]
    ({cos(\t)},{sin(\t)},0);
  \begin{pgfonlayer}{backedges}
    \draw[povmsphereback]
      plot[domain=120:300, samples=37, variable=\t, smooth]
      ({cos(\t)},{sin(\t)},0);
  \end{pgfonlayer}
}

% The panel radius in TikZ units.  Figure 2.2 draws at 2.6; a consumer that
% needs a different size renews this rather than copying the macros, so every
% panel everywhere stays one drawing (\renewcommand{\PanelScale}{2.47}).
\newcommand{\PanelScale}{2.6}

\newcommand{\PanelHeader}[2]{%
  \begin{scope}[xshift=#1mm, yshift=#2mm]
  \begin{scope}[scale=\PanelScale, tdplot_main_coords]%
  \Axes
  \Equator
}

% Variant for panels whose vertices don't touch the equator (just the tetrahedron):
% draws an outer silhouette circle to give the sphere a clear extent.
\newcommand{\PanelHeaderSil}[2]{%
  \begin{scope}[xshift=#1mm, yshift=#2mm]
  \draw[povmsphere] (0,0) circle (\PanelScale);
  \begin{scope}[scale=\PanelScale, tdplot_main_coords]%
  \Axes
  \Equator
}

\newcommand{\PanelFooter}[1]{% {panelTitle}
  \end{scope}% close inner tdplot scope
  \node[panelname] at (0, -2.85) {#1};
  \end{scope}% close outer xshift scope
}

